
\documentclass[../../main.tex]{subfiles}
\begin{document}

% 年度の大きな表示
\begin{center}
  {\fontsize{48pt}{58pt}\selectfont \textbf{1977}\normalsize\textbf{年度}}
\end{center}

\vspace{10mm}

% 本文

\vspace{15mm}

% 出題テーマと難易度の表
\noindent\textbf{出題テーマと難易度}

\vspace{5mm}

\begin{center}
  \shadowbox{%
    \begin{tabular}{|c|p{8cm}|c|}
      \hline
      \textbf{問題番号} & \textbf{テーマ} & \textbf{難易度} \\
      \hline
      第1問 & 円と接線 & \\
      \hline
      第2問 & 確率 & \\
      \hline
      第3問 & 接線のなす角 & \\
      \hline
      第4問 & 不等式の証明 & \\
      \hline
      第5問 & 積分方程式 & \\
      \hline
      第6問 & 数列と行列（選択問題） & \\
      \hline
    \end{tabular}%
  }
\end{center}

\vspace{15mm}

% 解答
\noindent\textbf{解答}

\vspace{5mm}

\begin{center}
  \shadowbox{%
    \renewcommand{\arraystretch}{1.6}
    \begin{tabular}{|c|c|p{8cm}|}
      \hline
      \textbf{問題番号} & \textbf{小問} & \textbf{答え} \\
      \hline
      第1問 & - & $Q\left(8r^4-8r^2+1,\ 4r(2r^2-1)\sqrt{1-r^2}\right)$ \\
      \hline
      第2問 & - & $\dfrac{27}{50}$ \\
      \hline
      第3問 & - & $P(1,1)$ \\
      \hline
      第4問 & - & 証明問題 \\
      \hline
      \multirow{2}{*}{第5問} & (1) & 証明問題 \\
      \cline{2-3}
                             & (2) & $f(x)=\sin x+\cos x,\ c=\pi$ \\
      \hline
      \multirow{2}{*}{第6問} & (a)-1 & 証明問題 \\
      \cline{2-3}
                             & (a)-2 & 証明問題 \\
      \hline
    \end{tabular}%
    \renewcommand{\arraystretch}{1.0}
  }
\end{center}

\vspace{15mm}

% 解答時間
\noindent\textbf{試験形態}

\vspace{5mm}

\begin{center}
  \shadowbox{%
    \begin{tabular}{p{8cm}}
      （要確認）
    \end{tabular}%
  }
\end{center}

\end{document}
